\scanpage{85}
\mspage{42}
d'ailleurs $G_m(S)$ n'a pas d'él. infiniment divisible\deleted{\unsure{s}}, donc
\[
(22)\quad G_m(S)\longrightarrow G_m(S)^\wedge\simeq H^1(S,T)
\]
\margin{(cas ``arithm.'')}
est injectif. Ainsi, quand $\unsure{\gamma_i|q_i}=0$ i.e. quand (13) admet des scindages ``géométriques'' \inserted{(\unsure{car} il suffit pour cela qu'il existe une section tt court)} ceux-ci forment un torseur sous $\deleted{\unsure{H^1(S,T)}}G_m(S)^{\unsure{\wedge}}$, qui s'identifie à un sous-torseur du $(H^1(S,T)=G_m(S)^\wedge)$-torseur de tous les scindages de (13). Pour que la ``description profinie de la géom. alg. absolue sur $\mathbb Q$'' soit complète, il y faudrait également caractériser (en termes de \inserted{cette} description profinie) ce sous-ensemble remarquable.

Je voudrais enfin comprendre encore comment une section d'extensions du type (1) peut se ``spécialiser'' en une section du type (2), dans le cas de \unsure{corps de} valuation. Pour ceci, je reprends
\margin{84}
